\subsection{An exterior-matched native coherence and its limits}
\label{subsec:native-localized-clock}

A separate construction tests localization directly on the native binary
carrier. With endpoint weights $\theta=(25,25,8,\ldots,8)$, its faithful
reference is $\pi_{ij}=\theta_i\theta_j/7505$. Choose the actual seam
$A=\{1,7\}$. For each $k\notin A$, replace the reference block on
$\{|1,k\rangle,|7,k\rangle\}$ by
\[
 \rho_k=\frac{\theta_k}{7505}
 \begin{pmatrix}25&8\\8&8\end{pmatrix},
\]
leaving every other entry unchanged. Each determinant is
$136(\theta_k/7505)^2>0$; normalization and the diagonal are unchanged.
The full exterior partial trace of $\rho-\pi$ vanishes. The global half
trace distance is $D_0=776/7505$.

Every complete native record acts by a partial injection. Equal records
preserve the signed difference of two load configurations, so a surviving
$|1,k\rangle\langle7,k|$ becomes $|1,l\rangle\langle7,l|$.
Its exterior partial trace is zero. Induction therefore proves equality
of the joint exterior and entire native record tape for $\rho$ and $\pi$
at every finite history length. The preparation archive is distinct and
is retained separately. This is a specified mixed state, not an arbitrary
logical-qubit encoding. Its coherent core cannot move under these branches.

For the uniform all-seam scheduler, the ten spectator amplitudes evolve
through an exact reversible substochastic matrix $Q$. With $n$ the number
of native attempts, their contrast obeys
\[
 D_n=\sum_k v_n(k),\qquad
 \sum_{n\ge0}\frac{D_n}{D_0}
 =\frac{\theta^\mathsf T(I-Q)^{-1}\mathbf1}
        {\sum_{k\notin A}\theta_k}.
\]
The integrated relative visibility is approximately $6.73354934$ attempts
for C, exactly $10/3$ for E, and approximately $4.45400705$ for the equal
classical mixture. For E, $D_n/D_0=(7/10)^n$. These are decays in native
attempt count. They are not particle lifetimes in a calibrated physical time.

The reference flow supplies phase $\exp[-is\log(25/8)]$, where $s$ is
modular parameter. A supplied Ramsey analyzer followed by a complete
native-current read has
\[
 \mathbb E j=-\frac{1649}{15010}+D_n\cos[s\log(25/8)].
\]
Despite that phase, the same state satisfies
$\operatorname{Tr}[(\rho-\pi)(-\log\pi)]=0$ and
$D(\rho\Vert\pi)=S(\pi)-S(\rho)>0$. Clock sensitivity and an information
deficit do not imply a nonzero energy or mass. Compiled rotations prepare
all 66 mixture columns on the original 512 source addresses; the replay
retains all 128 raw histories of a requested C-then-E continuation.
The exact scheduler calculation separately covers all 30 seams.

Exact arbitrary-qubit correction has a stronger restriction. The existing
full-78-state flagged-code theorem excludes double occupations and every
native-edge pair from any exact code of logical dimension at least two,
and makes each compressed occupation scalar. For any two-port core,
exterior events whose every core extension is excluded have code
probability zero but positive reference probability. The full core census
therefore gives exterior half trace distance at least $96/395$.
This applies to arbitrary exact subspace codes, without assuming a sharp
core occupation. For $N_H=n_0+n_1$ on the permitted support, write
$V^*N_HV=\nu I$. Up to scalar phase,
\[
 V^*\pi^{-is}V=[1-\nu+\nu e^{-is\log(25/8)}]I,
 \qquad
 p_{\rm leak}=4\nu(1-\nu)\sin^2\!\bigl(s\log(25/8)/2\bigr).
\]
An invariant exact code has $\nu=0$ or $1$ and a scalar logical clock.
Approximate codes, other references and additional controlled services are
outside this exclusion. The native construction supplies no protected
boundary charge. Its source, exact rational checks and complete histories
are retained in \texttt{model/native\_localization}.

\subsection{Moving the full band with a counted clock service}
\label{subsec:full-band-counted-clock}

The native two-level clock cannot extend to the full Weyl/contact band by
merely adding independent reference copies. Put $q=25/8$. A single native
reference has eigenvalues $(64/7505)q^m$, with $m=0,1,2$ and multiplicities
$45,20,1$. Every finite product has modular gaps in $\mathbb Z\log q$.
The eta calibration fixes $\beta U=\log q$. At total momentum zero, the
constant singlet obeys
\[
 \langle d,Td\rangle=\langle d,T^3d\rangle=0,
 \qquad \langle d,T^2d\rangle=6c^2=3U^2/32.
\]
The compression to $d,Td/\|Td\|$ has matrix
$U\bigl(\begin{smallmatrix}1&\sqrt{3/32}\\\sqrt{3/32}&0\end{smallmatrix}\bigr)$.
Its characteristic polynomial at $13/12$ is $-1/288$; min--max and
$\|T\|<U/2$ give $13/12<E_{\rm pair}/U<3/2$.
The eta dark state has energy zero. This noninteger difference excludes
an invariant isometry intertwining both field levels with any finite
independent native product at the same calibration. More generally,
$\|(A-aI)V-VH/U\|\le\varepsilon$ for an integer-spectrum native generator
$A$, any scalar $a$ and isometry $V$ implies $\varepsilon>1/24$.
Each of the two field energies must lie within $\varepsilon$ of the shifted
integer spectrum, whereas their separation is more than $1/12$ from any
integer. This is an inherited-reference exclusion, not a prohibition on
another reference or a counted simulation service.

Such a service is executed on the complete 8,128-dimensional sector using
all 64 band momenta, scattering states and Weyl tastes. The additional
reference premise is the globally normalized faithful density
$\tau=e^{-\beta H}/\operatorname{Tr}e^{-\beta H}$. It gives
$\tau^{is}A\tau^{-is}=e^{-i\beta sH}Ae^{i\beta sH}$ on the entire algebra.
Separate normalization of momentum fibers fails on cross-fiber operators.
The reference is supplied mathematically; its native selection and source
preparation are not established.

The native record counts infer $\widehat\beta$, and a declared source-control
service uses model duration $t=-\widehat\beta s$. Existing compiled local
hopping/contact words and counted fermionic routes implement a symmetric
32-step split. All local occupation inputs are checked; the lifted count
is 26,238,982 source pulses, rather than a literal replay of that many
pulses on a global processor. The count acquisition separately uses
31,457,280 native attempts. A moving band packet is entangled with an
untouched qubit, and all four terminal occupation-pointer outputs survive.
For $|s|\le1$, the confidence event gives
$|\widehat\beta-\beta|\le3/80$ and $|t|\le5/4$. Thus the complete average
half-diamond comparison obeys
\[
 \varepsilon\le\frac{9|s|}{256}
 +\frac{(7/4)^3(5|s|/4)^3}{3\,32^2}
 +\frac{381169}{10^{12}}<0.038565.
\]
The bound is uniform over all sector inputs and finite untouched references;
rejected count records give the same retained stop on both sides. No success
probability divides the bound. It compares ideal completed source words;
floating contraction residuals and hardware errors are separate diagnostics
and implementation premises. Packet preparation, coherent loading and
export, analogue control and the common field reference are declared inputs.
The source-bound execution in \texttt{model/band\_clock} therefore joins a
moving finite band to a record-derived rate under an explicit added premise;
it does not identify native repair with physical Hamiltonian motion.

\paragraph{Clock budget and energy offsets.}
The full-field enclosure
$\|H-(35/96)I\|\le77/96$ improves the rate-error term to
$77|s|/2560$. With $n$ symmetric split steps the bound becomes
\[
 \varepsilon_n\le\frac{77|s|}{2560}
 +\frac{(7/4)^3(5|s|/4)^3}{3n^2}
 +\frac{381169}{10^{12}}.
\]
At $|s|=1$, $n=21$ gives a bound below $0.037991$; $n=20$
exceeds the original 32-step ceiling. Thus 21 is minimal within this bound
and ceiling, not among every possible simulation. The complete retained
source schedule uses 17,219,334 pulses, saving 9,019,648. It retains all
8,128 states and four terminal outcomes. Acquisition, preparation and the
global-reference premises are unchanged; hardware and floating errors are
outside this ideal source bound.

The finite energy audit separates further identification questions. If
the generator inside each momentum fiber is known up to a scalar, a
connected tree of 63 admitted coherent comparisons would determine all
64 offsets up to one common scalar. These are mathematical access conditions,
not constructed native detectors. Even such access cannot detect $H\mapsto
H+cN$ within the fixed $N=2$ carrier: that change is $2cI$ and leaves its
normalized Gibbs state and dynamics unchanged. The eta-pair/dark contrast
$U$ remains an internal energy difference. Vacuum-relative rest energy
requires a further sector comparison or independently established preparation
law; dispersion curvature concerns a different kinematic contract.

For the native localized coherence, complete raw histories preserve equality
of its diagonals with the reference. If $R_n$ and $T_n$ are their respective
pushforwards under the same complete instrument, then
$\operatorname{Tr}[(R_n-T_n)(-\log T_n)]=0$ on the image support at every
finite history length. This uses no stationary state on the growing tape.
It extends the zero-energy boundary to the full archive; positive relative
entropy and decay in native attempt count do not furnish a mass.

\paragraph{Source-only conservation under classical feedback.}
For the faithful native $H_{66}$ reference $\pi$ and a conditional
Heisenberg attempt map $\Phi(X)=\sum_bK_b^*XK_b$, completeness and reference
invariance give
\[
 \operatorname{Re}\langle X,(I-\Phi)X\rangle_\pi
 =\tfrac12\sum_b\|XK_b-K_bX\|_\pi^2,
 \qquad \langle X,Y\rangle_\pi=\operatorname{Tr}(\pi X^*Y).
\]
Positive all-seam exploration and the complete branch algebra therefore
make every fixed source observable scalar. Applied separately at each
classical controller state, this excludes nonconstant source-only
conservation on all admitted source/controller inputs; it does not assume
a stationary growing archive. The exact weighted comparison gives
$\Gamma\ge g_0/10240>9.3468\,10^{-7}$ per native attempt throughout the
declared bounded family, where $g_0=\beta_{20}/20$ uses the retained exact
twenty-step population minorization. Hence
$\|X-\operatorname{Tr}(\pi X)I\|_\pi\le
\|X-\Phi X\|_\pi/\Gamma$.
Exogenous varying selectors sharing $\pi$ also contract; correlated
record-feedback mixing is not inferred. Joint relational charges,
restricted preparations, active quantum recovery and the nonfaithful
$H_{78}$ stationary support are outside this source-only result.
